%
%  chapter5.tex — Implementation and Current State
%
\chapter{Implementation and Current State}
\label{cha:implementation}

\section{Development Status}
\label{sec:dev_status}

Play4Change has been implemented across six development phases. Table~\ref{tab:dev_status} summarizes the current state of each major component.

\begin{table}[ht]
    \caption{Development status by component (April 2026).}
    \label{tab:dev_status}
    \vskip0.5em
    \centering
    \begin{tabular}{llp{5.5cm}}
    \toprule
    Component & Status & Notes \\
    \midrule
    Shared Kotlin Module     & \checkmark Done & API contracts, serialization, domain value objects. \\
    Backend — Identity BC    & \checkmark Done & Magic link + Google/Facebook OAuth, JWT, token family revocation, auth hardening. \\
    Backend — Content BC     & \checkmark Done & Topic creation, async generation, PDFBox + Jsoup extraction, MinIO file storage, task versioning. \\
    Backend — Learning BC    & \checkmark Done & Enrollments, day-index progression, struggle detection, adaptive branching, roadmap assembly. \\
    Backend — Gamification   & In Progress & Streaks and points implemented; badges and leaderboards pending. \\
    Backend — Peer Review BC & \checkmark Done & Reviewer selection, 3-verdict majority-vote finalization, social feedback, reviewer rewards. \\
    AI Pipeline (fixed path) & \checkmark Done & LangChain4j + Mistral, pgvector deduplication, batch @Scheduled, static fallback. \\
    AI Pipeline (adaptive)   & \checkmark Done & Similarity-based reuse (pgvector), StruggleAnalysisPrompt, error pattern classification. \\
    Observability            & \checkmark Done & Micrometer + Prometheus + Grafana, 7 custom domain metrics, provisioned dashboard. \\
    Redis Caching            & \checkmark Done & CachePort with silent fallback, 4 cache domains with TTLs. \\
    KMP Client (Android)     & In Progress & Onboarding, auth, topic browser, daily challenge flow, roadmap implemented. Gamification UI pending. \\
    KMP Client (iOS)         & In Progress & Shared components compile; SwiftUI integration in progress. \\
    Landing Page + Admin     & In Progress & Marketing content live; admin panel for topic management in progress. \\
    CI/CD                    & Configured & Lint, build, unit tests active; integration tests wired. \\
    Docker Compose           & \checkmark Done & Full stack (server, PostgreSQL+pgvector, Redis, MinIO, Prometheus, Grafana). \\
    \bottomrule
    \end{tabular}
\end{table}

\section{Key Implementation Decisions}
\label{sec:impl_decisions}

\subsection{Type Safety Across the System Boundary}
\label{ssec:typesafety_impl}

The shared Kotlin module is the most impactful single engineering decision in the project. All HTTP request and response bodies are defined as \texttt{@Serializable} Kotlin data classes in \texttt{commonMain}. Both the Ktor HTTP client and the Spring Boot controller use these same classes. A change to an API response that is not reflected in the shared module produces a compile-time error, not a runtime crash.

\subsection{Arrow Either — Railway Oriented Programming}
\label{ssec:arrow_impl}

The project began with a custom \texttt{Result<T, E>} sealed class (\texttt{Result.Success}, \texttt{Result.Failure}, \texttt{ResultScope}, \texttt{resultWrapper}) built intentionally from scratch to internalize the Railway Oriented Programming pattern. After comprehension was established, the codebase was migrated to Arrow \texttt{Either<AppError, T>} for production reasons: active maintenance, industry recognition, and rich combinators (\texttt{mapError}, \texttt{recover}, \texttt{zip}). The migration was mechanical (Table~\ref{tab:arrow_migration}). The \texttt{AppError} sealed hierarchy is domain-specific and unchanged. See ADR-001 in Appendix~\ref{app:adrs}.

\begin{table}[ht]
    \caption{Custom Result → Arrow Either migration map.}
    \label{tab:arrow_migration}
    \vskip0.5em
    \centering
    \begin{tabular}{ll}
    \toprule
    Custom & Arrow equivalent \\
    \midrule
    \texttt{Result<T, E>} & \texttt{Either<E, T>} \\
    \texttt{AppResult<T>} & \texttt{Either<AppError, T>} \\
    \texttt{Result.Success(data)} & \texttt{Either.Right(data)} \\
    \texttt{Result.Failure(error)} & \texttt{Either.Left(error)} \\
    \texttt{resultWrapper \{ \}} & \texttt{either \{ \}} \\
    \texttt{.bind()} & \texttt{.bind()} (same API) \\
    \texttt{ensure(cond) \{ err \}} & \texttt{ensure(cond) \{ err \}} (same) \\
    \texttt{.fold(onSuccess, onFailure)} & \texttt{.fold(onLeft, onRight)} \\
    \bottomrule
    \end{tabular}
\end{table}

\subsection{Passwordless Auth Hardening}
\label{ssec:auth_hardening}

A post-implementation security audit identified 11 gaps between the ADR-011 design intent and the initial implementation. Key remediations (ADR-016):

\begin{itemize}
    \item \textbf{G1}: Magic link raw tokens were stored in plaintext. Fixed: only \texttt{SHA-256(token)} is stored; the raw token travels only in the email link and the single verify request.
    \item \textbf{G2}: The token claim check was not atomic (TOCTOU race condition). Fixed: a single \texttt{UPDATE \dots WHERE used = false RETURNING *} CTE atomically claims the token.
    \item \textbf{G3}: Google OAuth audience validation was silently skipped when \texttt{GOOGLE\_CLIENT\_ID} was blank. Fixed: blank client ID is a hard startup error; audience is always validated.
    \item \textbf{G5}: Refresh rotation discarded the email claim. Fixed: the \texttt{User} entity is loaded by PK during refresh to re-embed the email.
    \item \textbf{G6}: Logout invalidated only the submitted token. Fixed: the entire token family is revoked on logout.
    \item \textbf{G7}: Public routes returned 401 when a stale Bearer token was present. Fixed: \texttt{JwtAuthFilter} continues the chain on parse failure for public routes.
\end{itemize}

\subsection{pgvector for Content Deduplication and Struggle Reuse}
\label{ssec:pgvector_impl}

PostgreSQL with the \textbf{pgvector} extension serves dual roles: task deduplication in the nightly batch and similarity-based reuse in adaptive branch generation. The cosine distance operator (\texttt{<=>}) is accessed via \texttt{JdbcTemplate} (not JPQL, which does not support custom operators). A custom \texttt{VectorConverter} handles bidirectional conversion between Kotlin \texttt{FloatArray} and pgvector's string format. Standard CRUD uses JPA; vector queries use \texttt{JdbcTemplate} — the right tool for each operation. The decision to use pgvector over a dedicated vector database (Qdrant, Weaviate) is documented in ADR-007: at current content volumes, pgvector's performance is equivalent, and it eliminates a deployment component and an additional network hop.

\subsection{Peer Review — Three-Verdict Majority}
\label{ssec:peer_review_impl}

The final task of each topic (\texttt{TODO\_ACTION}) requires the learner to photograph a real-world application of their learning. Three enrolled peers review the submission. Two verdicts (CORRECT or INCORRECT) trigger finalization synchronously on the third verdict — no async background job needed since finalization is bounded to one assignment update and at most four enrollment point updates. Reviewer selection uses application-level prioritization (unreviewed first, then partial, then random within tier) rather than \texttt{ORDER BY RANDOM()} to avoid full-table scans. A unique database constraint (\texttt{submission\_assignment\_id, reviewer\_user\_id}) provides a secondary safety net against duplicate assignments. After submitting a verdict, reviewers see the percentage breakdown of all verdicts — social feedback without prior anchoring bias (ADR-014).

\subsection{Adaptive Day Progression}
\label{ssec:day_impl}

Day index is computed from the \texttt{X-Timezone} HTTP header using \texttt{ChronoUnit.DAYS.between(LocalDate, LocalDate)} — calendar days, not seconds. The \texttt{DayIndexCalculator} is a pure Kotlin object: zero dependencies, zero mocking required in tests. The adaptive roadmap inserts struggle-branch nodes after the current day's task and resumes the fixed path after branch completion — no retry gate, no conditional advancement.

\section{Encountered Challenges}
\label{sec:challenges}

\subsection{PDFBox 3.x Breaking API Change}
The PDF text extraction library changed its entry point from \texttt{PDDocument.load()} to \texttt{Loader.loadPDF()} in version 3.x. This was discovered at compile time (correct) rather than runtime (the previous behaviour with a version mismatch).

\subsection{Mistral Prompt Engineering}
Generating consistently structured JSON from an LLM required iterative prompt engineering. The solution combined few-shot examples in the prompt, explicit JSON schema in the system message, defensive markdown-fence stripping in the parser (\texttt{removePrefix("```json")}), and \texttt{runCatching} per element so malformed tasks skip without aborting the batch.

\subsection{KMP/iOS Memory Model}
Decompose component callbacks and \texttt{StateFlow} collectors must be managed through the Decompose lifecycle to prevent memory leaks on Kotlin/Native (iOS). The final integration pattern uses \texttt{MutableValue} (Decompose-native) rather than \texttt{StateFlow}, rendering synchronously on the first composition without a \texttt{LaunchedEffect} frame delay.

\subsection{Known Gaps (Documented in ADR-007)}
Two known implementation gaps are tracked:
\begin{itemize}
    \item Fresh adaptive branch embeddings are not stored after generation, preventing the reuse library from growing from the most expensive generation path. Fix requires adding \texttt{struggleEventId} to \texttt{StruggleContext} and calling \texttt{storeStruggleEmbedding()} after branch persistence.
    \item \texttt{runBlocking} inside \texttt{@Transactional} in \texttt{CourseService} blocks a thread for the duration of Mistral inference while holding a database connection. Fix requires separating the save and generation transactions.
\end{itemize}

\section{Project Repository Structure}
\label{sec:project_structure}

The project is organized as a monorepo:

\begin{itemize}
    \item \texttt{/shared} — Shared Kotlin module (KMP): API contracts, value objects, validation.
    \item \texttt{/mobile} — KMP mobile application (\texttt{androidApp}, \texttt{iosApp}, \texttt{commonMain}).
    \item \texttt{/server} — Spring Boot backend.
    \item \texttt{/ai-agent/api} — \texttt{TaskGenerationPort} interface and models.
    \item \texttt{/ai-agent/langchain} — LangChain4j + Mistral implementation.
    \item \texttt{/landing} — TypeScript/React landing page and admin panel.
    \item \texttt{/infra} — Docker Compose, Prometheus config, Grafana dashboards, CI/CD.
    \item \texttt{/docs/adr} — Architecture Decision Records (also compiled in Appendix~\ref{app:adrs}).
\end{itemize}
